\clearpage
\section{Lean verification of Part I}\label{app:lean}
The source package \nolinkurl{lean/} proves selected Part~I results in Lean~4:
Theorem~3.3(a)--(c) and the elementary-amenability part of (d) for the
generated permutation group $\Gamma$, Theorem~2.3 conditionally, and the
printed, partial or existential forms listed below. It is the Lean core of
Part~I, not a formalization of the whole paper.

\subsection*{What is checked}
\nolinkurl{Statements.lean} freezes the selected propositions,
\nolinkurl{Proofs.lean} proves them and \nolinkurl{Audit.lean} checks that
they match the original declarations below. A row credits only the form
stated in it. In Lean, $\Gamma$ is the subgroup of
$\operatorname{Perm}(V\times S_3)$ generated by the five displayed
permutations $p,q,t,c,a$. The first column uses this paper's numbering;
Theorem~3.3 is the detailed Theorem~A. Some Lean names and docstrings keep
older numbering.

\begin{center}\footnotesize
\begin{minipage}{\linewidth}
\setlength{\tabcolsep}{4pt}
\begin{tabular}{@{}>{\raggedright\arraybackslash}p{.15\linewidth}>{\raggedright\arraybackslash}p{.30\linewidth}>{\raggedright\arraybackslash}p{.49\linewidth}@{}}
\toprule Printed result & Lean name & Checked form\\\midrule
Theorem 3.3(a) & \nolinkurl{Lamps.theoremA_a_unconditional} & The printed profile lower-bound shape, for $\Gamma$.\\[4pt]
Theorem 3.3(b) & \nolinkurl{Lamps.theoremA_b}; \nolinkurl{Lamps.theoremA_b_profile} & The explicit model upper bound $2^{500\cdot24^2 r\log_2(2+r)}$ for $r\ge2$, and the corresponding profile upper bound.\\[4pt]
Theorem 3.3(c) & \nolinkurl{Lamps.theoremA_c_unconditional} & The printed unitary lower-bound shape with existential positive constants: the scale is $N=\lfloor\log_2\lfloor X\rfloor\rfloor$, and $e_\Delta$ is existential (from little-o asymptotics), not explicit.\\[4pt]
Theorem 3.3(d) & \texttt{Lamps.}$\Gamma$\texttt{\_elementaryAmenable} & Partial: elementary amenability only.\\[4pt]
Theorem 3.3 & \nolinkurl{Lamps.theoremA} & Partial: (a), (b) for the chunk $E$, (c) and the elementary-amenability part of (d), for $\Gamma$; finite presentation is omitted.\\[4pt]
Theorem 2.1 & \nolinkurl{Lamps.counting_criterion_final} & Partial: the permutation/model/profile conclusion only; not the unitary criterion.\\[4pt]
Corollary 2.2 & \nolinkurl{Lamps.cyclic_criterion_final} & Partial: the permutation/profile conclusion only; not the unitary conclusion.\\[4pt]
Theorem 2.3 & \nolinkurl{Lamps.sweep_criterion} & Conditional on \nolinkurl{Lamps.PermSweepBound} (below): the whole-suffix permutation form; no elementary-letter variant.\\
\bottomrule
\end{tabular}\end{minipage}\end{center}

\begin{center}\footnotesize
\begin{minipage}{\linewidth}
\setlength{\tabcolsep}{4pt}
\begin{tabular}{@{}>{\raggedright\arraybackslash}p{.15\linewidth}>{\raggedright\arraybackslash}p{.30\linewidth}>{\raggedright\arraybackslash}p{.49\linewidth}@{}}
\toprule Printed result & Lean name & Checked form\\\midrule
Lemma 3.2 & \nolinkurl{Lamps.rho_rel} & The printed 25-word certificate.\\[4pt]
Lemma 4.1 & \nolinkurl{Lamps.exchange_sem} & Conditional: a semantic transporter estimate with supplied support, stabilizer-word and Dehn-bound data; not the full short-transporter specialization.\\[4pt]
Lemma 4.2 & \nolinkurl{Lamps.localIdsFamily} & Existential: the abstract family \texttt{LocalIdsFamily 2000 700}, with a Dehn bound at $2000n$ (the paper uses $1500n$), common area $10D+700n$ and existential translations and transvections; not the separate area-$1$ and ${<}40n+1$ estimates.\\[4pt]
Theorem 5.1 & \nolinkurl{Lamps.pairAreaStatement} & Existential: transport words with a polynomial pair-area bound and an unspecified constant (internal bound $49N^2L$; the paper states $20N^2L$); not their prescribed translation evaluations, $|T_u|\le21N\max(3,N)$ or $C_0=45000$.\\[4pt]
Proposition 6.3 & \nolinkurl{Lamps.prop_upper}; \nolinkurl{Lamps.prop_upper_profile} & Model and profile bounds for every fixed chunk of word length $\le\ell$, with $\log_2(2+r)$.\\[4pt]
Lemma B.1 & \nolinkurl{Lamps.unitaryS3Round} & Partial: unitary $x,c$ with $x^2=c^3=1$, $xcx=c^*$, $\|x-A\|_2\le e$ and $\|c-AB\|_2\le9e$; not the $8e$ and $10e$ estimates.\\[4pt]
Theorem B.3 & \nolinkurl{Lamps.unitaryPhaseRobust} & Partial: a weaker consequence after substituting pairwise commutator error bounds and dropping the positive part; not the displayed bound in the actual ordered-product errors.\\
\bottomrule
\end{tabular}\end{minipage}\end{center}

\noindent\textbf{Conditional premise.} The proposition
\nolinkurl{Lamps.PermSweepBound} is not an axiom and is not proved here; it is
defined in
\begin{quote}\small
\nolinkurl{lean/vendor/finite-lean/FiniteLean/Hypotheses/PermSweepBound.lean}.
\end{quote}
For every finite
permutation domain and every family of exact local $S_3$ actions whose
three-cycles move at least a quarter of the points, with commutator error at
most $1/20000$ against every whole ordered suffix and each of the six normal
forms, it asserts $\log_2|\Omega|\ge N/16$: the entropy step of the sweep
proof. Only Theorem~2.3 and its numbered wrapper take it; Theorem~3.3 does
not use it.

\noindent\textbf{Conventions:} finite profiles use normalized Hamming distance
with $\varepsilon=1/r$ and positive model sizes; unitary bounds use
$(\operatorname{Tr}X^*X/D)^{1/2}$, and an empty dimension type makes the
separation hypothesis unsatisfiable. The vendored entropy results are finite
dimensional, use natural logarithms and keep the positive-definiteness
hypotheses of their types.

\noindent\textbf{Not formalized:} the identification of $\Gamma$ with the full
wreath product $(\bigoplus_V S_3)\rtimes A$ (generation is Lemma~3.1(b)); the
exact $\mathbb Z^2$ quotient and kernel description of Lemma~3.1(c) (the
amenability proof uses a locally finite-by-abelian description instead);
finite presentation (Proposition~7.1 and Theorem~3.3(d));
Propositions~7.2--7.3; the ladder results of Section~8; the
\nolinkurl{PermSweepBound} entropy step and the elementary-letter $1/60000$
sweep variant; the unitary conclusions of Theorem~2.1 and Corollary~2.2 as selected statements (the source proves the unitary counting criterion as \nolinkurl{Lamps.counting_criterion_unitary}, relative to the proved interfaces \nolinkurl{UnitaryS3Round}, \nolinkurl{UnitaryPhaseRobust} and \nolinkurl{UnitaryMargin}, and uses it for Theorem~3.3(c); it is not a selected root)
and the general unitary sector statement of Lemma~B.2; Lemmas~6.1--6.2 and
A.1--A.2, Proposition~A.3 and Lemma~C.4 as numbered statements; Parts~II--III
and Appendices~C--I. \nolinkurl{lean/evidence/PAPER-COVERAGE.tsv} lists every
numbered result.

\subsection*{Trust base}
The proofs use Lean~4.35.0 and the Mathlib commit
\begin{quote}\small
\nolinkurl{56a005855ea17e0b0958f73ec801ae140b0a1679};
\end{quote}
all pins are in \nolinkurl{lean/}. Every declaration in the table and every
numbered wrapper depends only on Lean's three standard axioms
\texttt{propext}, \texttt{Classical.choice} and \texttt{Quot.sound};
\nolinkurl{lean/AXIOMS.txt} is the exact axiom output. The proofs pass two
independent kernel checkers, \texttt{leanchecker} and \texttt{nanoda}.
\nolinkurl{lean/evidence/KERNEL-REPLAY.json} records the kernel replay (export digest, checker binaries, exits) and the run of the reproduce commands on 6 November 2026; \nolinkurl{lean/evidence/CERTIFIED-SOURCE-HASHES.json} pins the checked sources.

\subsection*{Reproduce}
With the pinned toolchain on PATH, Git available for fetching the pinned
dependencies, and one compiler job at a time, run from \texttt{lean/}:
\begin{verbatim}
export LEAN_NUM_THREADS=1
lake exe cache get
lake build
lake env lean -j1 -M6144 Audit.lean
\end{verbatim}
\texttt{Audit.lean} prints the type of each checked declaration followed by its
axioms; the axiom lines match \nolinkurl{lean/AXIOMS.txt}.
